% Rúbrica: Procesamiento de datos. Fuente: notebooks/00_exploracion.ipynb, bloques 1 y 2, §3.5, y
% backend/analisis/limpieza.py.
\section{Procesamiento de datos}\label{sec:procesamiento}

El procesamiento tiene dos destinos. El modelo de riesgo se entrena con data-v1 tal como se
publicó, porque solo usa variables del juego y no lee texto. La limpieza produce otro conjunto,
el limpio, que usan el análisis de texto (sección~\ref{sec:exploracion-texto}) y la medición de
los motivos (sección~\ref{sec:interpretabilidad-motivos}).

\subsection{Validación}\label{sec:procesamiento-validacion}

Antes de todo, el notebook 00 corre los chequeos de la tabla~\ref{tab:calidad}
(sección~\ref{sec:datos-calidad}). Los de rango que nunca deben fallar, como un minuto negativo
o un voto distinto de 0 y 1, son aserciones\footnote{Una aserción es una comprobación dentro del
código: si no se cumple, la ejecución se detiene con un error.}, así que el notebook no sigue si
uno falla. También se revisan las copias que cruzan la partición: \cifraTextosEntreJuegos{}
textos idénticos, de \cifraPalabrasDeUnaCopia{} palabras o más, aparecen en más de un juego, y
\cifraCopiasEnFoldsDistintos{} de ellos quedan en folds distintos. Un modelo de texto podría
verlos en entrenamiento y en prueba, y por eso la limpieza quita los duplicados. La partición
misma está congelada en un archivo (sección~\ref{sec:modelacion-validacion}).

\subsection{Limpieza: marcar, no borrar}\label{sec:procesamiento-limpieza}

Cada regla es una función de \texttt{backend/analisis/limpieza.py}, y se aplican en el orden de
la tabla~\ref{tab:limpieza}. Solo dos quitan filas: los textos idénticos, de los que se queda la
reseña más antigua, y las reseñas vacías. Los duplicados quitan \cifraDuplicadosQuitados{}
filas y las vacías \cifraVaciasQuitadas{}, y quedan \cifraResenasLimpias{} reseñas. Las demás reglas agregan
una columna que marca la reseña y la dejan en su lugar. Las cortas, por ejemplo, son el
\cifraCortasPorResena{} de las reseñas, y quitarlas cambiaría qué reseñas se estudian, no solo
cuánto ruido hay. La normalización pasa el texto a minúsculas, quita el BBCode\footnote{El
formato que Steam permite en las reseñas, con etiquetas como \texttt{[h1]} o
\texttt{[spoiler]}.} y las direcciones web, y abre las contracciones: «don't» queda como «do
not», y la negación, como palabra suelta.

\begin{table}[htbp]
  \centering\small
  % Generado por documento/generar_figuras.py. No se edita a mano.
\begin{tabularx}{\textwidth}{>{\raggedright\arraybackslash}X >{\raggedright\arraybackslash}p{4.4cm} rr}
\toprule
Regla, en orden & Qué se hace & Afectadas & Quedan \\
\midrule
Texto idéntico de 8 palabras o más & se quitan; queda la más antigua & 483 (0.4\,\%) & 123,489 \\
Reseña vacía & se quitan & 545 (0.4\,\%) & 122,944 \\
Menos de 3 palabras & se marcan & 27,944 (22.7\,\%) & 122,944 \\
Otro idioma, según lingua & se marcan & 1,637 (1.3\,\%) & 122,944 \\
Plantilla de casillas & se marcan & 191 (0.2\,\%) & 122,944 \\
Normalización del texto & se agrega \texttt{texto\_norm} & 94,303 (76.7\,\%) & 122,944 \\
\bottomrule
\end{tabularx}

  \caption{Bitácora de la limpieza sobre data-v1. El porcentaje es sobre las reseñas que llegan a
  cada regla; en la normalización, las afectadas son las que cambian de texto. Fuente:
  \texttt{limpiar} de \texttt{backend/analisis/limpieza.py}.}
  \label{tab:limpieza}
\end{table}

La limpieza no se aplica al modelo de riesgo. Se midió qué le pasaría con la misma partición
congelada, para que la diferencia saliera de las filas y no del reparto de los juegos. Sin los
duplicados, su PR-AUC pasa de \cifraPRAUCModelo{} a \cifraPRAUCSinDuplicados{}, y sin las
vacías, a \cifraPRAUCSinVacias{}. Sin las cortas subiría a \cifraPRAUCSinCortas{}, pero sobre
otra población y con más prevalencia, así que los dos números no se comparan. El modelo se queda
con data-v1 tal cual. El conjunto limpio se guarda con un sha256 de su contenido, y quien lo
rehaga puede comprobar que obtuvo el mismo.

\subsection{Idioma}\label{sec:procesamiento-idioma}

La ingesta pidió reseñas en inglés (\texttt{language=english}), y Steam entrega las que llevan
esa etiqueta. Para revisarla se usó lingua\footnote{Una biblioteca de detección de idioma. Solo
decide cuando tiene confianza; los textos muy cortos o sin sentido quedan como
indeterminados.}: el \cifraIdiomaInglesPorResena{} de las reseñas limpias está en inglés o no se
pudo determinar, y \cifraOtroIdioma{} están escritas en otro idioma, sobre todo ruso, español y
portugués. Como el filtro lo aplicó Steam, esa cifra mide qué tan bien etiquetaron sus reseñas
los autores, no qué idiomas se usan en la tienda. Las reseñas en otro idioma se marcan y se
quedan. Al modelo de riesgo no le afectan, porque no lee texto, pero a los motivos sí, porque sus
palabras clave están en inglés (sección~\ref{sec:interpretabilidad-motivos}).

\subsection{Precio imputado}\label{sec:procesamiento-precio}

Los \cifraPagoSinPrecioEntrenamiento{} juegos de pago que llegaron sin precio
(tabla~\ref{tab:calidad}) entran al modelo con precio 0: la variable es el logaritmo de uno más
el precio, y los faltantes se llenan con 0. El modelo los lee como si costaran 0, aunque no los
marca como gratis, y como el coeficiente del precio es positivo, eso tiende a bajar su riesgo
estimado. La ficha de esos juegos lo avisa (sección~\ref{sec:interpretabilidad-avisos}).

\subsection{Variables que no se conocen al comprar}\label{sec:procesamiento-variables}

Una variable solo puede entrar al modelo si se conoce antes de comprar. Si no, el modelo
aprendería de algo que todavía no existe cuando la persona decide, y eso es una
fuga\footnote{Hay fuga cuando un modelo usa información que no tendría en el momento de
predecir. Sus métricas salen mejores de lo que serán en uso.}. La tabla~\ref{tab:variables}
repasa las candidatas. Las del juego vienen de la tienda y entran. Las del autor de la reseña
no: casi todas existen solo porque la reseña ya se escribió, y la biblioteca del autor, que sí
podría declararse antes, se probó y su aporte cupo en el ruido entre folds
(sección~\ref{sec:modelacion-bmas}). En el sitio, el perfil se declara en un formulario y sirve
para la afinidad: no mueve el riesgo.

\begin{table}[htbp]
  \centering\small
  % Generado por documento/generar_figuras.py. No se edita a mano.
\begin{tabularx}{\textwidth}{>{\raggedright\arraybackslash}X >{\raggedright\arraybackslash}p{4.2cm} >{\raggedright\arraybackslash}p{4.6cm}}
\toprule
Variable & Cuándo se conoce & ¿Entra al modelo de riesgo? \\
\midrule
\texttt{es\_gratis}, precio (en logaritmo) y descuento & en la tienda, antes de comprar & sí \\
Si tiene nota de Metacritic, y la nota & en la tienda, antes de comprar & sí; la nota que falta se rellena con la mediana, junto a la bandera \\
\texttt{num\_games\_owned}: juegos del autor & en la reseña; en el sitio se declara «compras al año» & no: su aporte cabe en el ruido entre folds, y 0 es un perfil privado \\
\texttt{num\_reviews}: reseñas del autor & al descargar; cambia con el tiempo & no: no está disponible al momento de la compra \\
\texttt{steam\_purchase}, \texttt{received\_for\_free} y \texttt{written\_during\_early\_access} & en la reseña & no: no están disponibles al momento de la compra \\
\texttt{playtime\_at\_review} y \texttt{voted\_up} & en la reseña & no: definen $Y$ \\
Texto de la reseña & en la reseña & no: sería una fuga, porque $Y$ sale de esas mismas reseñas \\
\bottomrule
\end{tabularx}

  \caption{Variables candidatas: cuándo se conocen y si entran al modelo de riesgo. Fuente:
  \texttt{construir\_features} de \texttt{backend/modelado/entrenar\_baseline.py} y la sección
  «Variables del jugador» del notebook 00.}
  \label{tab:variables}
\end{table}
